\begin{itemize}
    \item \textbf{Gegenstand:} Erzeugung realistischer klinischer Daten für das Heidelberger Lehre-KIS durch die Kombination von generativen Modellen, regelbasierten Systemen und Sensordatenintegration.
    \item \textbf{Bedeutung:} Steigerung der Lehrqualität durch authentische Fallbeispiele, die angehende Fachkräfte auf die Komplexität und Dynamik rechnerbasierter Anwendungssysteme vorbereiten.
    \item \textbf{Problematik:} Bestehende synthetische Datensätze sind oft statisch und bilden die zeitliche Variabilität klinischer Parameter unzureichend ab. Zudem erschweren proprietäre Schnittstellen die Einbindung realer Medizinprodukte in Lehrumgebungen.
    \item \textbf{Motivation:} Das Ziel ist die Schaffung eines „lebendigen“ KIS, das durch die Fusion von algorithmisch erzeugten Krankheitsverläufen und echten physiologischen Datenströmen (Wearables/Sensoren) eine hohe Plausibilität der bereitgestellten Daten erreicht. Durch diesen technologieoffenen Ansatz können sowohl moderne Sprachmodelle als auch bewährte Simulationsmethoden genutzt werden, um die Lücke zwischen statischer Theorie und dynamischer klinischer Praxis zu schließen.
\end{itemize}